\documentclass[11pt,a4paper]{article}

\usepackage[utf8]{inputenc}
\usepackage[T1]{fontenc}
\usepackage[brazilian]{babel}
\usepackage{lmodern}
\usepackage{microtype}
\usepackage{geometry}
\usepackage{xcolor}
\usepackage{graphicx}
\usepackage{tabularx}
\usepackage{booktabs}
\usepackage{array}
\usepackage{enumitem}
\usepackage{amsmath}
\usepackage{hyperref}

\geometry{margin=2.2cm}
\setlength{\parindent}{0pt}
\setlength{\parskip}{0.55em}
\setlist{nosep,leftmargin=1.4em}
\renewcommand{\arraystretch}{1.18}

\definecolor{navy}{HTML}{17324D}
\definecolor{blue}{HTML}{246B9E}
\definecolor{lightblue}{HTML}{EAF4FA}
\definecolor{orange}{HTML}{E87524}
\definecolor{lightorange}{HTML}{FFF1E7}
\definecolor{green}{HTML}{277A53}
\definecolor{lightgreen}{HTML}{EAF6EF}
\definecolor{graytext}{HTML}{4C5964}

\hypersetup{
  colorlinks=true,
  linkcolor=blue,
  urlcolor=blue,
  pdftitle={Metodologia de Calibração e Rastreamento da Bola},
  pdfauthor={Projeto Robot Soccer Vision}
}

\newcommand{\conceptbox}[2]{%
  \begin{center}
  \fcolorbox{#1}{#1!8}{\parbox{0.91\linewidth}{#2}}
  \end{center}
}

\newcommand{\stage}[2]{%
  \fcolorbox{blue}{lightblue}{\parbox[c][1.18cm][c]{0.17\linewidth}{\centering\textbf{#1}\\[-1mm]\small #2}}
}

\begin{document}
\pagenumbering{gobble}

\begin{titlepage}
  \centering
  \vspace*{2.2cm}
  {\color{navy}\rule{\textwidth}{1.5pt}}\\[1.2cm]
  {\Huge\bfseries\color{navy} Metodologia de Calibração\\[0.25cm]
  e Rastreamento da Bola\par}
  \vspace{0.7cm}
  {\Large\color{graytext} Visão computacional para futebol de robôs VSSS\par}
  \vspace{1.5cm}
  \conceptbox{orange}{
    \textbf{Ideia central.} Um modelo robusto, porém caro, aprende com o vídeo e produz
    referências confiáveis. Depois, um modelo leve é escolhido para se aproximar dessas
    referências com baixa latência. Quadros difíceis são separados para revisão humana.
  }
  \vfill
  {\large Projeto Robot Soccer Vision\par}
  \vspace{0.2cm}
  {\color{graytext}Documento conceitual --- setembro de 2026\par}
  \vspace{1.2cm}
  {\color{navy}\rule{\textwidth}{1.5pt}}
\end{titlepage}

\pagenumbering{roman}
\tableofcontents
\newpage
\pagenumbering{arabic}

\section{O problema em poucas palavras}

A câmera observa o campo por cima e permanece fixa. A bola é laranja, aproximadamente
circular e quase não muda de tamanho na imagem. Essas três informações são muito úteis:

\begin{itemize}
  \item a \textbf{cor} reduz a busca às regiões parecidas com laranja;
  \item a \textbf{forma circular} rejeita linhas, etiquetas e manchas;
  \item o \textbf{raio quase constante} elimina objetos grandes ou pequenos demais.
\end{itemize}

O desafio é que a aparência da cor muda com iluminação, exposição, sombras e reflexos.
Além disso, o rastreador precisa responder rapidamente. Por isso, o sistema separa duas
funções:

\begin{center}
\begin{tabularx}{0.94\linewidth}{>{\bfseries\color{navy}}p{2.3cm}X}
\toprule
Modelo robusto & Usa várias técnicas e busca consenso. É mais lento e serve como
``professor'' durante calibração e avaliação.\\
Modelo leve & Usa apenas a combinação escolhida pelas ablações. É o modelo executado
em tempo real durante o jogo.\\
\bottomrule
\end{tabularx}
\end{center}

\conceptbox{green}{
  O modelo robusto não é uma verdade perfeita. Ele é uma \textbf{pseudo-referência}.
  Uma pequena amostra rotulada manualmente continua necessária para detectar erros que
  os dois modelos possam cometer juntos.
}

\section{Calibração inicial}

\subsection{Correção de perspectiva com um quadrilátero}

O usuário marca os quatro cantos do campo na ordem: superior esquerdo, superior direito,
inferior direito e inferior esquerdo. Esses pontos definem um quadrilátero. Uma
\textbf{homografia} transforma esse quadrilátero em um retângulo, removendo a inclinação
aparente causada por uma câmera lateralmente deslocada ou não perfeitamente perpendicular.

O quadrilátero é usado \textbf{somente para corrigir o enviesamento da imagem}. Ele não
converte pixels em milímetros e não inventa margens fora da região selecionada. O mapa da
transformação é calculado uma vez e reutilizado nos demais quadros.

\subsection{Amostra inicial da bola}

Na imagem corrigida, o usuário posiciona uma circunferência sobre a bola e ajusta seu raio.
Os pixels da parte interna dessa circunferência fornecem a primeira amostra de cor. Usa-se a
mediana dos valores, porque ela é menos sensível a um pequeno reflexo branco ou a um pixel
com ruído.

\begin{center}
\stage{1}{Quad do campo}\hspace{0.15em}$\longrightarrow$\hspace{0.15em}
\stage{2}{Raio da bola}\hspace{0.15em}$\longrightarrow$\hspace{0.15em}
\stage{3}{Cor inicial}\hspace{0.15em}$\longrightarrow$\hspace{0.15em}
\stage{4}{Busca automática}
\end{center}

\section{Representação de cor e adaptação à luz}

A imagem é convertida de BGR para \textbf{HSV}. Nesse espaço, cada pixel é descrito por:

\begin{itemize}
  \item \textbf{H (matiz):} o tipo de cor, por exemplo laranja;
  \item \textbf{S (saturação):} quão viva ou acinzentada é a cor;
  \item \textbf{V (valor):} o brilho.
\end{itemize}

Um pixel entra na máscara de bola quando sua distância para a cor aprendida está dentro das
tolerâncias de H, S e V. A matiz é circular: valores próximos das extremidades da escala
também podem representar cores vizinhas.

Depois de uma detecção de alta confiança, o sistema observa novamente a parte interna da
bola e atualiza a cor por uma \textbf{média móvel exponencial}:

\[
  \text{cor nova} = (1-\alpha)\,\text{cor anterior}
  + \alpha\,\text{cor observada}.
\]

O fator $\alpha$ é a taxa de aprendizagem. Um valor pequeno acompanha mudanças graduais
sem reagir demais a um único quadro. A atualização possui travas: exige boa forma circular,
boa confiança, saturação mínima e limita a mudança de matiz por quadro. Se a bola não for
encontrada, o modelo de cor fica congelado para evitar que o fundo seja aprendido por engano.

\section{Pré-processamento comparado}

O afinamento testa quatro alternativas. Nenhuma é sempre melhor; a escolha depende do
vídeo e do custo medido no computador real.

\begin{center}
\begin{tabularx}{\linewidth}{>{\bfseries}p{3.5cm}X}
\toprule
Alternativa & Intuição e principal compromisso\\
\midrule
Sem filtro & Preserva a máscara original e tem custo mínimo, mas é mais sensível a pixels isolados.\\
Gaussiano $5\times5$ & Suaviza variações locais antes da decisão, porém pode borrar a borda da bola.\\
Abertura circular & Erosão seguida de dilatação remove pequenos pontos, com risco de reduzir objetos muito pequenos.\\
Abertura + fechamento & Além de remover pontos, preenche pequenos buracos causados por reflexo, ao custo de mais operações.\\
\bottomrule
\end{tabularx}
\end{center}

Os elementos morfológicos são elípticos, não quadrados. Isso respeita melhor a geometria da
bola e evita introduzir preferência artificial por direções horizontal e vertical.

\section{Técnicas de detecção}

\subsection{Filtro casado disco-menos-anel}

O detector leve principal usa uma convolução com um núcleo circular pré-calculado. O centro
do núcleo possui pesos positivos e representa a bola. Um anel externo possui pesos
negativos e representa o fundo esperado ao redor dela.

\conceptbox{orange}{
  Uma bola laranja isolada produz resposta alta: o interior combina com o disco positivo e o
  entorno não ativa o anel negativo. Uma grande região laranja tende a invadir o anel e perde
  pontuação.
}

Como o raio muda pouco, o núcleo é criado uma vez. A convolução procura o ponto de maior
resposta. Essa operação é densa, simples e rápida em uma pequena região de busca.

\subsection{Componentes conexos}

A máscara também pode ser dividida em regiões conectadas. Para cada região, calculam-se
área, perímetro, círculo envolvente, circularidade e diferença para o raio esperado. Regiões
com tamanho incorreto são descartadas. Essa abordagem é barata quando a máscara contém
poucos objetos, mas pode sofrer quando muitas manchas têm a mesma cor.

\subsection{Transformada de Hough}

O modelo robusto usa ainda a transformada de Hough para círculos. Ela procura evidências de
bordas organizadas ao redor de um centro e de um raio. É mais cara, mas fornece uma fonte de
evidência diferente da cor preenchida e dos componentes conexos.

\section{O professor robusto}

O professor executa uma busca global e combina várias opiniões:

\begin{enumerate}
  \item filtro casado e componentes conexos;
  \item três raios próximos ao raio inicial;
  \item abertura e fechamento morfológicos;
  \item detecção circular por Hough;
  \item preferência temporal por posições compatíveis com o movimento recente.
\end{enumerate}

As propostas próximas são agrupadas. Cada grupo recebe votos e uma pontuação de acordo com
confiança, concordância espacial e proximidade da posição prevista. Só grupos com consenso
mínimo geram uma pseudo-etiqueta. A cor adaptativa do professor também só é atualizada após
uma decisão confiável.

Esse conjunto custa mais tempo que o permitido no jogo, mas reduz a chance de confiar em um
único tipo de erro. Ele é adequado para processar vídeos gravados fora do laço de controle.

\section{Autossupervisão e escolha de hiperparâmetros}

Autossupervisão significa criar um sinal de qualidade a partir do próprio vídeo, sem exigir
a posição medida por um sensor externo em todos os quadros.

O sistema experimenta variações da cor inicial, raio, tolerâncias HSV, tolerância de raio e
taxa de adaptação. Cada candidato é avaliado por quatro sinais:

\begin{itemize}
  \item \textbf{cobertura:} quantos quadros recebem uma detecção;
  \item \textbf{confiança:} quão forte é o consenso interno do professor;
  \item \textbf{ida e volta:} o vídeo processado do início ao fim deve concordar com o mesmo
  vídeo processado do fim ao início;
  \item \textbf{iluminação:} a posição deve permanecer semelhante após alterações artificiais
  de brilho e saturação.
\end{itemize}

O melhor professor produz pseudo-etiquetas para o vídeo completo. Em seguida são comparadas
16 combinações do aluno: quatro pré-processamentos, dois detectores e busca por região ligada
ou desligada. O critério final combina:

\[
  \text{objetivo} = \text{erro de centro normalizado}
  + \text{penalidade por falhas}
  + \text{penalidade por latência}.
\]

Esse procedimento seleciona um modelo que não é apenas rápido, nem apenas parecido com o
professor: ele procura um equilíbrio entre ambos.

\section{Modelo leve e controle de latência}

Durante o jogo, o modelo leve usa a velocidade recente para prever onde a bola estará e
recorta uma \textbf{região de interesse} (ROI). Cor, pré-processamento e detecção são
executados apenas nessa região. Se a bola for perdida, o sistema volta temporariamente para
uma busca global em meia resolução.

Outras medidas reduzem atraso:

\begin{itemize}
  \item o núcleo circular é pré-calculado;
  \item a captura mantém somente o quadro mais novo e descarta quadros antigos;
  \item o mapa de correção de perspectiva é pré-calculado;
  \item a interface informa tempo de processamento, idade do quadro e uso de ROI;
  \item a cor só se adapta depois de detecções confiáveis.
\end{itemize}

Latência deve ser medida da captura até o resultado. Uma taxa alta de quadros por segundo não
garante baixa latência se houver uma fila de imagens antigas.

\section{Aprendizado ativo e revisão humana}

Rotular todos os quadros é caro e pouco informativo, pois a maioria deles é fácil. O
aprendizado ativo calcula uma dificuldade para cada quadro usando:

\begin{itemize}
  \item baixa confiança e poucos votos do professor;
  \item diferença entre professor e aluno;
  \item movimento inesperado em relação à previsão temporal;
  \item pouca nitidez, medida pela variância do Laplaciano;
  \item baixo brilho e grande fração de sombras.
\end{itemize}

Forma-se um conjunto com os quadros mais difíceis e sorteiam-se $N$ exemplos desse conjunto.
O sorteio evita revisar apenas um trecho contínuo do vídeo. O usuário pode marcar o centro e
o raio ou pular quadros impossíveis por desfoque de movimento, escuridão ou oclusão. O motivo
do descarte também é guardado e pode revelar limitações da câmera ou da iluminação.

\section{Avaliação e interpretação correta}

Na avaliação, professor e aluno recebem exatamente o mesmo quadro corrigido. São registrados
erro entre centros, falhas do aluno, detecções exclusivas do aluno, confiança, latência e
quadros descartados pela captura. Os resultados agregados incluem média, mediana, percentil
95 e erro máximo, além de concordância relativa ao raio da bola.

\conceptbox{green}{
  \textbf{Como interpretar:} boa concordância com o professor mostra que o aluno preservou o
  comportamento do modelo robusto. Ela não prova que ambos estão corretos. A validação final
  deve comparar os dois com uma pequena amostra manual independente.
}

Para permitir um relatório científico posterior, cada etapa registra caminhos, hashes
SHA-256, parâmetros, versões das bibliotecas, revisão Git, resultados por quadro e artefatos
gerados. Assim, tabelas e gráficos podem ser reproduzidos sem depender de anotações feitas à
mão durante o experimento.

\section{Resumo visual}

\begin{center}
\fcolorbox{navy}{lightblue}{\parbox{0.92\linewidth}{
\centering
\textbf{Vídeo + calibração visual} $\rightarrow$
\textbf{professor robusto} $\rightarrow$
\textbf{pseudo-etiquetas}\\[0.35em]
$\downarrow$\\[-0.1em]
\textbf{ablações de baixo custo} $\rightarrow$
\textbf{modelo leve escolhido} $\rightarrow$
\textbf{avaliação e quadros difíceis}\\[0.35em]
$\downarrow$\\[-0.1em]
\textbf{revisão humana pequena} $\rightarrow$
\textbf{resultado confiável e documentado}
}}
\end{center}

\vfill
\begin{center}
\small\color{graytext}
Este documento descreve os princípios da metodologia implementada no projeto
Robot Soccer Vision. Ele prioriza intuição e interpretação, não detalhes de código.
\end{center}

\end{document}
